\documentclass[10pt,a4paper]{amsart}
\usepackage[margin=19mm]{geometry}
\usepackage{amsmath,amssymb,mathtools}
\usepackage[scr=rsfs,cal=euler]{mathalfa}
\usepackage{xfrac}
\usepackage[unicode,hidelinks,bookmarks=false]{hyperref}
\newcommand{\ie}{i.e.\ }
\newcommand{\cf}{cf.\ }
\newcommand{\resp}{respectively\ }
\newcommand{\Subsection}{\subsection}
\providecommand{\nobreakdash}{\nobreak\hbox{-}}
\setlength{\emergencystretch}{2em}
\multlinegap0pt
\newtheorem{thm}{Theorem}
\newtheorem{cor}{Corollary}
\title{$\ \ $ A realizability criterion for rational maps with three branch points}
\author{Zhiqiang Wei}
\date{Source: arXiv:2607.00818v2 (CC BY 4.0)}
\begin{document}
\maketitle

\noindent\emph{Source and licence.} $\ \ $ A realizability criterion for rational maps with three branch points. Version arXiv:2607.00818v2 (CC BY 4.0), distributed by arXiv under the Creative Commons Attribution 4.0 International licence, http://creativecommons.org/licenses/by/4.0/, as recorded in the arXiv licence metadata for this exact version and shown on the abstract page https://arxiv.org/abs/2607.00818v2. Full licence notice: \texttt{licenses/arxiv-2607.00818-CC-BY-4.0.txt}.\par\smallskip

\noindent\textit{Editorial scope.} Counted unit: the article's cor Co1 (source0.tex lines 322--380). The article's complete original proof of it is delivered with it (source0.tex lines 382--392, one \texttt{proof} environment of 11 lines). No mathematics is written, repaired or re-derived: the statement and the proof are the article's own lines, in the article's own notation, and no retained line is substituted or re-flowed.  Retained uncounted as the source-local context the proof needs: lines 278--314.  Nothing was omitted from any retained span, so the package carries no omission record.  No retained line contains a citation, so the wrapper carries no bibliography and no bibliography entry.  No retained line contains a figure environment, an image-inclusion command or drawing code, so no image is delivered and the package declares no asset dependency.  Editorial formatting only, in addition: an AMS \texttt{amsart} preamble that restates the article's own declarations and macros used by the retained lines, namely the environments \texttt{thm}, \texttt{cor} on separate counters, together with the standard packages those lines need.  The retained environments print in this document as Theorem 1, Corollary 1 in order of appearance, whereas the article prints the same items with the numbering of its own full document.  The complete native source of the article as distributed in the arXiv e-print for this exact version is bundled unchanged as \texttt{sources/204191/source0.tex}.
\begin{thm}\label{MThm2}
If \(\gamma_i \ge 2\) for each \(i=1,2,3\), then \(\mathcal D\) is realizable by a rational map \(f:\overline{\mathbb{C}}\to \overline{\mathbb{C}}\) with branch points \(0,\infty,1\), corresponding to \(\pi_0,\pi_\infty,\pi_1\), respectively, if and only if one of the following statements holds.

\begin{enumerate}
\item[(A)] The map \(f\) is of Type-A and one of the following conditions holds.
\begin{enumerate}
\item[(A-1)] For every permutation \(\sigma\) of \(\{1,2,3\}\), the set \(\{\widetilde{\pi}_0,\widetilde{\pi}_\infty,\widetilde{\pi}_1\}\) is realizable, where
\[
\widetilde{\pi}_0=[\underbrace{2,\ldots,2}_{2s+1}],~
\widetilde{\pi}_\infty=[\underbrace{3,\ldots,3}_{s+1},\underbrace{1,\ldots,1}_{s-1}],~
\widetilde{\pi}_1=[\gamma_{\sigma(1)}+3,\ \gamma_{\sigma(2)}+1,\ \gamma_{\sigma(3)}].
\]
\item[(A-2)] There exists a permutation \(\tau\) of \(\{1,2,3\}\) such that the set \(\{\widehat{\pi}_0,\widehat{\pi}_\infty,\widehat{\pi}_1\}\) is realizable, where
\[
\widehat{\pi}_0=[\underbrace{2,\ldots,2}_{2s+1}],~
\widehat{\pi}_\infty=[\underbrace{3,\ldots,3}_{s+1},\underbrace{1,\ldots,1}_{s-1}],~
\widehat{\pi}_1=[\gamma_{\tau(1)}+4,\ \gamma_{\tau(2)},\ \gamma_{\tau(3)}].
\]
\end{enumerate}

\item[(B)] The map \(f\) is of Type-B and one of the following conditions holds.
\begin{enumerate}
\item[(B-1)] There exists a permutation \(\tau\) of \(\{1,2,3\}\) such that the set \(\{\widetilde{\pi}_0,\widetilde{\pi}_\infty,\widetilde{\pi}_1\}\) is realizable, where
\[
\widetilde{\pi}_0=[\underbrace{2,\ldots,2}_{2s+1}],~
\widetilde{\pi}_\infty=[\underbrace{3,\ldots,3}_{s+1},\underbrace{1,\ldots,1}_{s-1}],~
\widetilde{\pi}_1=[\gamma_{\tau(1)}+3k,\ \gamma_{\tau(2)}+k,\ \gamma_{\tau(3)}].
\]
\item[(B-2)] There exists a permutation \(\sigma\) of \(\{1,2,3\}\) such that the set \(\{\widehat{\pi}_0,\widehat{\pi}_\infty,\widehat{\pi}_1\}\) is realizable, where
\[
\widehat{\pi}_0=[\underbrace{2,\ldots,2}_{2s+1}],~
\widehat{\pi}_\infty=[\underbrace{3,\ldots,3}_{s+1},\underbrace{1,\ldots,1}_{s-1}],~
\widehat{\pi}_1=[\gamma_{\sigma(1)}+4,\ \gamma_{\sigma(2)},\ \gamma_{\sigma(3)}].
\]
\end{enumerate}
\end{enumerate}
\end{thm}

\begin{cor}\label{Co1}
Let $k \ge 0$ be an integer. Then the following sets of candidate branch data are realizable:
\begin{enumerate}
\item $\displaystyle
\bigl\{
   \bigl[ \underbrace{2,\ldots,2}_{5+2k} \bigr],\
   \bigl[ \underbrace{3,\ldots,3}_{3+k},\underbrace{1,\ldots,1}_{1+k} \bigr],\
   [2+3k,\,3+k,\,5]
\bigr\}.
$
\item $\displaystyle
\bigl\{
   \bigl[ \underbrace{2,\ldots,2}_{5+2k} \bigr],\
   \bigl[ \underbrace{3,\ldots,3}_{3+k},\underbrace{1,\ldots,1}_{1+k} \bigr],\
   [2+3k,\,3,\,5+k]
\bigr\}.
$

\item $\displaystyle
\bigl\{
   \bigl[ \underbrace{2,\ldots,2}_{5+2k} \bigr],\
   \bigl[ \underbrace{3,\ldots,3}_{3+k},\underbrace{1,\ldots,1}_{1+k} \bigr],\
   [2,\,3+3k,\,5+k]
\bigr\}.
$

\item $\displaystyle
\bigl\{
   \bigl[ \underbrace{2,\ldots,2}_{5+2k} \bigr],\
   \bigl[ \underbrace{3,\ldots,3}_{3+k},\underbrace{1,\ldots,1}_{1+k} \bigr],\
   [2+k,\,3+3k,\,5]
\bigr\}.
$

\item $\displaystyle
\bigl\{
   \bigl[ \underbrace{2,\ldots,2}_{5+2k} \bigr],\
   \bigl[ \underbrace{3,\ldots,3}_{3+k},\underbrace{1,\ldots,1}_{1+k} \bigr],\
   [2+k,\,3,\,5+3k]
\bigr\}.
$

\item $\displaystyle
\bigl\{
   \bigl[ \underbrace{2,\ldots,2}_{5+2k} \bigr],\
   \bigl[ \underbrace{3,\ldots,3}_{3+k},\underbrace{1,\ldots,1}_{1+k} \bigr],\
   [2,\,3+k,\,5+3k]
\bigr\}.
$

\item $\displaystyle
\bigl\{
   \bigl[ \underbrace{2,\ldots,2}_{5+2k} \bigr],\
   \bigl[ \underbrace{3,\ldots,3}_{3+k},\underbrace{1,\ldots,1}_{1+k} \bigr],\
   [2,\,3,\,5+4k]
\bigr\}.
$
\end{enumerate}
\end{cor}

\begin{proof}
The proof of Theorem~\ref{MThm2} establishes the realizability of the base candidate branch datum
\[
\bigl\{
 [2,2,2,2,2],\
  [3,3,3,1],\
  [2,3,5]
\bigr\}.
\]
Items (1)--(6) then follow immediately from part (A) of Theorem~\ref{MThm2}. Item (7) follows from part (ii) of Theorem~\ref{MThm2}, as also shown in the proof of that theorem.
\end{proof}

\end{document}
